% conclusions.tex — Section 9: Conclusions

\section{Conclusions}\label{sec:conclusions}

We have shown that autonomous AI agents stress all nine behavioral invariants on which many banking fraud-detection systems rely. The resulting failures are not merely threshold-calibration problems where the monitored signal presupposes human embodiment: agents have no biological body to tire, no inherently persistent device fingerprint, no single embodied location, and near-zero marginal cost of identity creation. Invariants such as velocity bounds, behavioral biometrics, and KYC scarcity assumptions~\cite{bsa1970} therefore require reinterpretation when the transacting entity is software. Detection systems built only on these invariants risk confusing legitimate agent commerce with agent-orchestrated fraud because the human-era calibration target is no longer stable.

The five-signal detection framework introduced in this paper addresses this structural gap by detecting what is \emph{present} in agent behavior but \emph{absent} in human behavior: deterministic network topology, machine-precision value flows, superhuman economic rationality, programmatic temporal patterns, and cross-platform behavioral inconsistencies. Validated on 81,904 real on-chain USDC transactions from 665 ERC-8004-registered agents~\cite{erc8004} sourced via Dune Analytics~\cite{dune2026}, the framework achieves 95.4\% recall on the raw evaluation set (81.1\% after label cleaning) at a re-optimized threshold, though composite ROC-AUC is only 0.515---reflecting negative-class contamination rather than ranking power. At the fixed synthetic threshold, recall drops to 61.4\%, revealing a genuine distributional shift between synthetic and real data that threshold re-optimization partially masks. Seven of the eight attack chains defined in our threat taxonomy are detected at 100\% per-chain recall in injection experiments, including three scenarios outside current human-assumption detectors.

The eighth chain---swarm intelligence attacks involving coordinated multi-agent populations that are individually indistinguishable from human users---reveals the clearest architectural gap in the current framework. Per-address scoring cannot detect collective behavior that is designed to be individually invisible. Population-level detection operating at block-granularity time windows represents the next frontier, requiring new primitives for activation correlation, amount synchronization, and temporal clustering that remain future work in this study.

The transition this work anticipated is no longer purely prospective. The capability stack for agent-to-agent fraud---autonomous transaction execution, programmatic identity generation, and multi-agent coordination via platforms such as OpenClaw~\cite{openclaw2026}---has been operational since mid-2025, and by mid-2026 documented on-chain agent-payment incidents~\cite{grokbankr2026}, large-scale agent impersonation~\cite{datadome2026}, and an explicit industry pivot toward ``Know Your Agent'' monitoring~\cite{financialbrand2026,trmlabs2026} are matters of public record. What remains undocumented is systematic A2A-commerce fraud exploiting the specific invariant blind spots modeled here; a confirmed in-scope case would be the decisive test of the threat model. The detection framework, open-source reference implementation, and institutional recommendations presented in this paper are designed to let institutions measure and harden controls against the in-scope attack surface as it is exercised, rather than after exploitation is observed at scale.
